\documentclass[11pt]{article}
\usepackage[margin=1in]{geometry}
\usepackage{amsmath, amssymb}
\usepackage{parskip}
\usepackage{hyperref}
\usepackage{enumitem}

\newcommand{\Rt}{R(t)}
\newcommand{\lstar}{\ell^{*}}
\newcommand{\bhat}{\hat\beta}
\newcommand{\bbar}{\bar\beta}

\title{Does the global SVD test in F1.1 measure what we want?}
\author{}
\date{\today}

\begin{document}
\maketitle

\section{The question}

Aggregating logits across prompts before doing SVD might flatten prompt-specific structure into a common, easier-to-fit direction.
If so, the off-axis residual fraction \(\Rt\) reported in F1.1 could understate the failure of per-trajectory multiplicativity,
and the SVD-recovered \(\bbar(k,t)\) shown in F1.2-svd could be tracking the wrong axis entirely.

This note works through what global \(\Rt\) actually measures, where it can be fooled, and what the cheapest fix would be.

\section{The framework claim, written per-trajectory}

The imprint assumption is per-prompt:
\[
\ell_t^{(n,k)} \;=\; \beta_{\text{model}}^{(n,k,t)} \cdot \lstar^{(n)}_t,
\]
for trajectory \(n\), decoder index \(k\), step \(t\),
with \(\lstar^{(n)}_t\) the prompt-specific ``natural'' logit direction.
That is the claim we want to test.

\section{What F1.1 actually computes}

Global \(\Rt\) averages first and tests after:
\[
\bar\ell_t(k) \;=\; \frac{1}{N}\sum_{n=1}^{N} \beta_{\text{model}}^{(n,k,t)} \cdot \lstar^{(n)}_t,
\qquad
M_t \;=\; \bigl[\,\widetilde{\bar\ell}_t(1)\,\bigl|\,\dots\,\bigl|\,\widetilde{\bar\ell}_t(K)\,\bigr],
\]
\[
M_t \;=\; U\Sigma V^\top, \qquad \Rt \;=\; \frac{\sum_{i\geq 2}\sigma_i^2}{\sum_i \sigma_i^2}.
\]
\(\Rt=0\) means the \(K\) \emph{averaged} columns lie on a one-dimensional ray.
That is \emph{necessary} for per-trajectory multiplicativity but not \emph{sufficient}.
Two failure modes hide in the gap.

\section{Failure mode 1: zero-mean deviations cancel under averaging}

Suppose the real per-trajectory imprint has a quadratic kicker,
\[
\ell_t^{(n,k)} \;=\; \beta_k\,\lstar^{(n)}_t \;+\; \beta_k^2\,\eta^{(n)},
\qquad \eta^{(n)} \;\sim\; \text{i.i.d., zero-mean over prompts.}
\]
Per-prompt, the imprint is non-multiplicative.
But the aggregate kicker washes out,
\[
\bar\ell_t(k) \;=\; \beta_k \cdot \underbrace{\tfrac{1}{N}\!\sum_n \lstar^{(n)}_t}_{\mu_t}
\;+\; \beta_k^2 \cdot \underbrace{\tfrac{1}{N}\!\sum_n \eta^{(n)}}_{\to\, 0}
\;\approx\; \beta_k\cdot\mu_t.
\]
So \(M_t\) is rank-1, \(\Rt\approx 0\), and the global test \emph{certifies multiplicativity that isn't there.}
This is the case the plan's html already flagged:
\emph{``Path-dependent off-axis structure that is zero-mean across trajectories cancels in the averaged test.''}

\section{Failure mode 2: the principal direction picks up shared content}

The deeper concern.
Different prompts have different \(\lstar^{(n)}_t\), but they share statistical structure ---
common-token directions (``the'', ``of'', ``and'', \dots) get systematically larger logits across most natural English contexts.
Call that shared direction \(\mu_t^{\text{avg}}\).
When we average over \(N\) prompts,
\[
\bar\ell_t(k) \;\approx\; \beta_k \cdot \mu_t^{\text{avg}} \;+\; (\text{prompt-specific terms, smaller in } N).
\]
The SVD's principal left singular vector \(u_1\) lines up with \(\mu_t^{\text{avg}}\),
not with any single \(\lstar^{(n)}_t\).
Then
\[
\bbar(k,t) \;=\; \sigma_1 \cdot v_{1k}
\]
measures how much each \(\beta_k\)-column projects onto the \emph{average-natural} direction.
That projection can track \(\beta_k\) closely even if per-prompt imprints are doing weirder things,
as long as those weird things don't share a direction across prompts.

This is consistent with what F1.2-svd shows.
The SVD-recovered \(\log\bbar\) spans only \([-1.1,\,-0.2]\) at \(t=199\),
while the streaming \(\log\bhat\) (F1.2) spans \([-4.7,\,+1.1]\).
The SVD is finding a flatter, content-averaged direction;
the streaming estimator is responding to something prompt-specific.
They disagree because they are measuring different things, not because one is buggy.

\section{So what does \(\Rt\) measure?}

The \emph{consistent-across-prompts} off-axis structure.
The observed \(\Rt[199]\approx 0.155\) (F1.1) means that even after averaging out prompt-specific quirks,
roughly 15\% of the cross-decoder variance does not fit a single principal direction.
That is a real positive finding --- consistent non-multiplicativity exists.

But \(\Rt\) is a \emph{lower bound} on per-trajectory off-axis residual:
\begin{itemize}[leftmargin=2em]
\item \(\Rt > 0\) \(\Rightarrow\) consistent (across prompts) non-multiplicativity.
\item \(\Rt \approx 0\) \(\Rightarrow\) consistent with per-prompt multiplicativity \emph{or} with wildly non-multiplicative per-prompt imprints that happen to be zero-mean across prompts. Inconclusive.
\end{itemize}

\section{Why per-prompt R is too noisy to break the tie}

The plan anticipates this and defines a per-prompt SVD (F1.4).
For each prompt \(p\), average just that prompt's \(N=3\) runs, form \(M_t^{(p)} \in \mathbb{R}^{V\times K}\), and compute \(R^{(p)}(t)\).
With the imprint exactly multiplicative per-prompt, \(M_t^{(p)}\) is rank-1 and \(R^{(p)}(t) = 0\).

Sampling-noise floor for this estimator.
Model
\[
\ell_t^{(n,k)} \;=\; \beta_k\,\lstar^{(n)}_t \;+\; \varepsilon_t^{(n,k)},
\qquad \varepsilon_t^{(n,k)} \;\sim\; \text{i.i.d., per-entry variance } \sigma^2,
\]
and average over \(N\) runs at fixed prompt.
The K-column matrix \(M_t^{(p)}\) has signal Frobenius-squared \(\|\widetilde{\lstar}\|^2\,\|\beta\|^2\)
and noise Frobenius-squared \(V K \sigma^2 / N\).
SVD's top singular value captures the signal; the remaining \(K-1\) singular values capture the orthogonal noise.
Assuming \(\sigma \approx \sigma_\ell\)
(per-trajectory variance comparable to the natural logit magnitude, since both come from the same model on similar prefixes),
\[
\mathbb{E}\bigl[R^{(p)}\bigr]
\;\approx\;
\frac{(K-1)/N}{\|\beta\|^2 + K/N}.
\]
With \(K=6\), \(N=3\), and \(\|\beta\|^2 = \sum_k \beta_k^2 \approx 10\) (for the sweep \(\{0.5,0.7,1.0,1.3,1.7,2.0\}\)),
\[
\mathbb{E}\bigl[R^{(p)}\bigr]
\;\approx\;
\frac{5/3}{10 + 2}
\;\approx\;
0.14.
\]
The observed gap between per-prompt and global in F1.4 is \(\approx 0.32 - 0.16 = 0.16\),
so the bulk of the per-prompt curve is finite-\(N\) noise.
With \(N=3\) we cannot distinguish ``consistent + per-prompt noise'' from ``consistent + per-prompt failure.''

For the global estimator with \(N=96\) trajectories the same calculation gives
\[
\mathbb{E}\bigl[R^{\text{global}}\bigr]
\;\approx\;
\frac{5/96}{10 + 6/96}
\;\approx\;
0.005,
\]
so the global noise floor is \(\sim 0.5\%\) and the observed \(\Rt[199]\approx 0.155\) is genuine signal ---
but, again, only of the consistent-across-prompts component.

\section{What would actually test what we want}

\begin{enumerate}[leftmargin=2em]
\item \textbf{More runs per prompt.}
Bumping to \(N=30\) drops the per-prompt noise floor to \(\approx 0.02\) and makes F1.4 the direct test.
Cost: \(10\times\) GPU on the existing layout.

\item \textbf{Toy validation with controlled \(\lstar\).}
A single deterministic-context setting on actual transformer logits, so aggregation across prompts isn't a confound.
The walkthrough's analytic toy already does this for the corrector;
a logit-level toy would close the gap for the SVD test specifically.

\item \textbf{Decompose \(\Rt\) into shared vs.\ prompt-specific.}
Project each per-prompt \(u_1^{(p)}\) onto the global \(u_1\) and look at the distribution.
If the projections are all \(\approx 1\), the global SVD is finding the right direction.
If they spread, the global \(u_1\) is the wrong axis and the F1.2-svd compression is a pure aggregation artifact.
\end{enumerate}

\section{Why F2 is mostly unaffected}

F2.1 and F2.2 compare the streaming \(\bhat_t\) against the SVD-recovered \(\bbar_t\).
The streaming estimator is computed per-trajectory and only averaged \emph{after} taking logarithms,
so it is not subject to the directional-aggregation failure mode.
The gap shown in F2.1 (growing with \(t\), anti-symmetric in \(\log\beta_{\text{dec}}\)) is real evidence
that the model is not Bayes-optimal under the assumed prior,
independent of whatever the global SVD's \(\bbar\) is doing.
So even if F1.1 and F1.2-svd require caveats,
F2's headline ``calibration gap grows with \(t\)'' is solid.

\section{Bottom line}

F1.1 is honest about \emph{consistent-across-prompts} structure:
\(\Rt[199] \approx 0.155\) is genuine non-multiplicativity.
But it cannot rule out larger, prompt-specific failures,
and the SVD's recovered \(\bbar\) is biased toward a content-averaged direction
that does not reflect any single prompt's natural axis.
The per-prompt SVD (F1.4) is the right test,
but at \(N=3\) runs/prompt it is dominated by finite-sample noise.
Cheapest path to a definitive answer: rerun with \(N\geq 30\).

\end{document}
